\documentclass[10pt]{article}

\usepackage[margin=0.75in]{geometry}
\usepackage{amsmath}
\usepackage{amssymb}
\usepackage{booktabs}
\usepackage{longtable}
\usepackage{array}
\usepackage{microtype}
\usepackage{underscore}
\usepackage{enumitem}
\usepackage{hyperref}

\allowdisplaybreaks
\setlength{\tabcolsep}{4pt}
\renewcommand{\arraystretch}{1.15}

\title{Capped Flexible-Docking Ordered-Sortie MILP\\
\large Formulation, Written to Match the Implementation}
\author{}
\date{}

\begin{document}
\maketitle

\noindent This document writes out the mixed-integer linear program (MILP) that is
actually implemented in \texttt{model/capped\_flexible\_docking\_ordered\_sortie\_model.py}.
It is not a transcription of the original PDF formulation. Every symbol, set, and
constraint below corresponds to named variables, constraints, or helper blocks in the code, so the formulation can be checked against the source
directly. Where this model departs from the original PDF formulation, that is
called out explicitly in Section~\ref{sec:deviations}.

\section{Sets and Indices}
\label{sec:sets}

{\small
\begin{longtable}{@{}p{0.17\textwidth}p{0.79\textwidth}@{}}
\toprule
Symbol & Meaning \\
\midrule
$C = \{1,\dots,n\}$ & Customers. $n = |C|$. \\
$P = \{0, n+1\}$ & Depots: $0$ is the start depot, $n+1$ is the end depot. \\
$V$ & $V = \{0,1,\dots,n+1\} = C \cup P$: all nodes. \\
$A \subseteq V \times V$ & Directed arc set (from \texttt{arcs.csv}); $(i,j)\in A$ means a truck may travel directly from $i$ to $j$. \\
$T = \{1,\dots,|T|\}$ & Trucks, $|T| = \texttt{NUM\_TRUCKS}$. \\
$D = \{1,\dots,|D|\}$ & Drones, after lossless fleet trimming (Section~\ref{sec:trimming}). \\
$R = \{1,\dots,|R|\}$ & Robots, after lossless fleet trimming. \\
$S^d$ & Set of candidate ordered drone sorties (Section~\ref{sec:sorties}). \\
$S^r$ & Set of candidate ordered robot sorties. \\
\bottomrule
\end{longtable}
}

\section{Ordered Sortie Candidates}
\label{sec:sorties}

Before the MILP is built, \texttt{generate\_ordered\_sorties} enumerates, for every
arc $(i,k)\in A$ and every ordered subsequence of customers of length
$1,\dots,\texttt{MAX\_CUSTOMERS\_PER\_SORTIE}$, a candidate \emph{ordered sortie}
$s$ with:

\begin{itemize}
\item launch node $\iota(s) = i$, recovery node $\kappa(s) = k$;
\item ordered customer sequence $\sigma(s) = (\sigma_1(s),\dots,\sigma_{m(s)}(s))$, $\sigma_l(s)\in C$;
\item payload $q(s) = \sum_{c\,\in\,\sigma(s)} q_c$, kept only if $q(s) \le$ platform capacity;
\item travel distance $\delta(s) = \sum_{(u,v)\,\in\,\mathrm{edges}(s)} \mathrm{dist}_{uv}$ (drone: \texttt{distance\_matrix}; robot: \texttt{robot\_distance\_matrix}), kept only if $\delta(s) \le$ platform endurance;
\item duration $\tau(s) = \sum_{(u,v)\,\in\,\mathrm{edges}(s)} \mathrm{time}_{uv}$;
\item cost $\gamma(s) = \sum_{(u,v)\,\in\,\mathrm{edges}(s)} \mathrm{cost}_{uv}$,
\end{itemize}

\noindent where
\[
\mathrm{edges}(s) = \big((\iota(s),\sigma_1(s)),(\sigma_1(s),\sigma_2(s)),\dots,(\sigma_{m(s)}(s),\kappa(s))\big).
\]
For each $(i,k)\in A$, only the
\texttt{TOP\_SORTIES\_PER\_TRUCK\_LEG} candidates with the most customers (ties
broken by lowest cost, then duration) are kept. $S^d$ uses drone time/distance
matrices and $(Q_d, E_d)$; $S^r$ uses robot time/distance matrices and $(Q_r,
E_r)$.

Therefore, an ``integer optimal solution'' from the solver is optimal over the
retained capped candidate-sortie sets $S^d$ and $S^r$. It is not a proof of
global optimality over every possible ordered sortie in the unrestricted problem.

\section{Parameters}
\label{sec:parameters}

{\small
\begin{longtable}{@{}p{0.19\textwidth}p{0.78\textwidth}@{}}
\toprule
Symbol & Meaning (source) \\
\midrule
$q_i$ & Demand of customer $i\in C$ (\texttt{customers.csv}). \\
$e_i,\, \bar{e}_i$ & Time-window open/close time of customer $i$ (\texttt{open\_time}, \texttt{close\_time}). \\
$\theta_i$ & Service time of customer $i$ (\texttt{service\_time}). \\
$\mathrm{dist}_{ij}$ & Truck/drone travel distance, $i,j\in V$ (\texttt{distance\_matrix}). \\
$\mathrm{tt}_{ij}$ & Truck travel time, $i,j\in V$ (\texttt{truck\_time\_matrix}). \\
$\mathrm{dt}_{ij}$ & Drone travel time, $i,j\in V$ (\texttt{drone\_time\_matrix}). \\
$\mathrm{rdist}_{ij}$ & Robot travel distance, $i,j\in V$ (\texttt{robot\_distance\_matrix}). \\
$\mathrm{rt}_{ij}$ & Robot travel time, $i,j\in V$ (\texttt{robot\_time\_matrix}). \\
$Q_t, Q_d, Q_r$ & Truck / drone / robot payload capacity. \\
$E_d, E_r$ & Drone / robot endurance, as a maximum route distance per sortie. \\
$T_{\max}$ & Route-duration limit (10 h in the TDRP-TW data). \\
$C_w, C_{veh}$ & Truck time-cost and distance-cost coefficients. \\
$C^d_w, C_{drone}$ & Drone time-cost and distance-cost coefficients (\texttt{C\_w\_drone}, \texttt{C\_drone}); these are explicit in the processed cases, with code fallbacks only for backward compatibility. \\
$C^r_w, C_{rob}$ & Robot time-cost and distance-cost coefficients (\texttt{C\_w\_r}, \texttt{C\_rob}). \\
$f_t, f_d, f_r$ & Fixed activation cost per used truck / selected drone sortie / selected robot sortie (\texttt{truck\_fixed\_cost}, \texttt{drone\_fixed\_cost}, \texttt{robot\_fixed\_cost}; values $30, 10, 8$ from Malik et al., VRP-DR, arXiv:2505.23584, Table 3). \\
$\lambda_T, \lambda_W$ & Penalty weights for route-duration excess and time-window lateness. \\
$\mu^{dep}_d$ & Drones a used truck may launch from the depot (\texttt{DRONES\_CARRIED\_AT\_DEPOT}). \\
$\mu^{trk}_d$ & Distinct drones a single truck may use (\texttt{MAX\_DRONES\_PER\_TRUCK}). \\
$\mu^{dep}_r$ & Robots a used truck may launch from the depot (\texttt{ROBOTS\_CARRIED\_AT\_DEPOT}). \\
$\mu^{trk}_r$ & Distinct robots a single truck may use (\texttt{MAX\_ROBOTS\_PER\_TRUCK}). \\
$M$ & Big-M, $M = \max(T_{\max},\ 10\cdot|V|\cdot\max(\mathrm{tt},\mathrm{dt},\mathrm{rt}),\ 10000)$. \\
\bottomrule
\end{longtable}
}

\subsection{Travel-cost coefficients}

\begin{align}
\mathrm{tcost}_{ij} &= \mathrm{tt}_{ij}\,C_w + \mathrm{dist}_{ij}\,C_{veh}, & (i,j)\in A \label{eq:tcost}\\
\mathrm{dcost}_{ij} &= \mathrm{dt}_{ij}\,C^d_w + \mathrm{dist}_{ij}\,C_{drone}, & i,j\in V \label{eq:dcost}\\
\mathrm{rcost}_{ij} &= \mathrm{rt}_{ij}\,C^r_w + \mathrm{rdist}_{ij}\,C_{rob}, & i,j\in V \label{eq:rcost}
\end{align}

\noindent Sortie cost $\gamma(s)$ in Section~\ref{sec:sorties} sums $\mathrm{dcost}$
(drone sorties, $s\in S^d$) or $\mathrm{rcost}$ (robot sorties, $s\in S^r$) over
$\mathrm{edges}(s)$.

\section{Decision Variables}

{\small
\begin{longtable}{@{}p{0.13\textwidth}p{0.17\textwidth}p{0.66\textwidth}@{}}
\toprule
Symbol & Code name & Meaning \\
\midrule
$x_{v,i,j}\in\{0,1\}$ & \texttt{x[v,i,j]} & Truck $v$ traverses arc $(i,j)\in A$. \\
$\eta_v\in\{0,1\}$ & \texttt{used\_truck[v]} & Truck $v$ is used. \\
$\mathrm{vis}_{v,i}\in\{0,1\}$ & \texttt{visit[v,i]} & Truck $v$ visits customer $i\in C$. \\
$u_{v,i}\in[0,n]$ & \texttt{u[v,i]} & MTZ rank of customer $i$ on truck $v$'s route. \\
$y_{d,s}\in\{0,1\}$ & \texttt{y[d,sid]} & Drone $d$ flies ordered sortie $s\in S^d$. \\
$h^L_{v,d,s}\in\{0,1\}$ & \texttt{h\_launch[v,d,sid]} & Truck $v$ launches drone $d$ for sortie $s$. \\
$h^R_{v,d,s}\in\{0,1\}$ & \texttt{h\_recover[v,d,sid]} & Truck $v$ recovers drone $d$ for sortie $s$. \\
$\xi_{v,d}\in\{0,1\}$ & \texttt{drone\_truck[v,d]} & Drone $d$ interacts with truck $v$ (launch and/or recovery). \\
$z_{r,s}\in\{0,1\}$ & \texttt{z[r,sid]} & Robot $r$ performs ordered sortie $s\in S^r$. \\
$g^L_{v,r,s}\in\{0,1\}$ & \texttt{g\_launch[v,r,sid]} & Truck $v$ launches robot $r$ for sortie $s$. \\
$g^R_{v,r,s}\in\{0,1\}$ & \texttt{g\_recover[v,r,sid]} & Truck $v$ recovers robot $r$ for sortie $s$. \\
$\zeta_{v,r}\in\{0,1\}$ & \texttt{robot\_truck[v,r]} & Robot $r$ interacts with truck $v$. \\
$a^t_{v,i}\ge 0$ & \texttt{a\_t[v,i]} & Arrival time of truck $v$ at node $i\in V$. \\
$a^d_i\ge 0$ & \texttt{a\_d[i]} & Drone arrival time at node $i\in V$. \\
$a^r_i\ge 0$ & \texttt{a\_r[i]} & Robot arrival time at node $i\in V$. \\
$\ell^t_i\ge 0$ & \texttt{late\_t[i]} & Truck time-window lateness at customer $i\in C$. \\
$\ell^d_i\ge 0$ & \texttt{late\_d[i]} & Drone time-window lateness at customer $i\in C$. \\
$\ell^r_i\ge 0$ & \texttt{late\_r[i]} & Robot time-window lateness at customer $i\in C$. \\
$\rho^t_v\ge 0$ & \texttt{route\_late\_t[v]} & Route-duration excess of truck $v$ over $T_{\max}$. \\
$\rho^*\ge 0$ & \texttt{route\_late\_global} & Global route-duration excess (objective uses this). \\
\bottomrule
\end{longtable}
}

\subsection{Derived service indicators}

\begin{align}
\sigma^t_i &= \sum_{v\in T} \mathrm{vis}_{v,i}, & i\in C \label{eq:truck-service}\\
\sigma^d_j &= \sum_{d\in D}\sum_{\substack{s\in S^d\\ j\,\in\,\sigma(s)}} y_{d,s}, & j\in C \label{eq:drone-service}\\
\sigma^r_j &= \sum_{r\in R}\sum_{\substack{s\in S^r\\ j\,\in\,\sigma(s)}} z_{r,s}, & j\in C \label{eq:robot-service}
\end{align}

\noindent\emph{Code:} \texttt{truck\_service}, \texttt{drone\_service}, \texttt{robot\_service}.

\section{Objective Function}

\subsection{Variable (distance/time) cost}
\begin{align}
\mathrm{VarCost}^t &= \sum_{v\in T}\sum_{(i,j)\in A} \mathrm{tcost}_{ij}\,x_{v,i,j} \label{eq:varcost-truck}\\
\mathrm{VarCost}^d &= \sum_{d\in D}\sum_{s\in S^d} \gamma(s)\,y_{d,s} \label{eq:varcost-drone}\\
\mathrm{VarCost}^r &= \sum_{r\in R}\sum_{s\in S^r} \gamma(s)\,z_{r,s} \label{eq:varcost-robot}\\
\mathrm{VarCost} &= \mathrm{VarCost}^t + \mathrm{VarCost}^d + \mathrm{VarCost}^r \label{eq:varcost-total}
\end{align}

\subsection{Fixed activation cost}
\begin{align}
\mathrm{FixCost}^t &= f_t \sum_{v\in T} \eta_v \label{eq:fixcost-truck}\\
\mathrm{FixCost}^d &= f_d \sum_{d\in D}\sum_{s\in S^d} y_{d,s} \label{eq:fixcost-drone}\\
\mathrm{FixCost}^r &= f_r \sum_{r\in R}\sum_{s\in S^r} z_{r,s} \label{eq:fixcost-robot}\\
\mathrm{FixCost} &= \mathrm{FixCost}^t + \mathrm{FixCost}^d + \mathrm{FixCost}^r \label{eq:fixcost-total}
\end{align}

\subsection{Operating cost and penalties}
\begin{align}
\mathrm{OpCost} &= \mathrm{VarCost} + \mathrm{FixCost} \label{eq:opcost}\\
\mathrm{Pen}_T &= \lambda_T\, \rho^* \label{eq:pen-route}\\
\mathrm{Pen}_W &= \lambda_W \left(\sum_{i\in C}\ell^t_i + \sum_{i\in C}\ell^d_i + \sum_{i\in C}\ell^r_i\right) \label{eq:pen-tw}\\
\mathrm{TotalPen} &= \mathrm{Pen}_T + \mathrm{Pen}_W \label{eq:totalpen}
\end{align}

\subsection{Objective}
\begin{align}
\min\quad Z = \mathrm{OpCost} + \mathrm{TotalPen} \label{eq:objective}
\end{align}

\noindent\emph{Code:} the objective and KPI expression block. Capacity and endurance are hard constraints
(Sections~\ref{sec:capacity}--\ref{sec:endurance}) and carry no penalty term; only
route-duration excess and time-window lateness are penalized.

\section{Constraints}

\subsection{Customer service (assignment)}
\begin{align}
\sigma^t_i + \sigma^d_i + \sigma^r_i = 1, \qquad i\in C \label{eq:service}
\end{align}
\noindent\emph{Code:} \texttt{customer\_service\_\{i\}}. Every customer is
served exactly once, by exactly one of truck, drone, or robot.

\subsection{Capacity (hard)}
\label{sec:capacity}
\begin{align}
\sum_{i\in C} q_i\, \mathrm{vis}_{v,i} &\le Q_t, & v\in T \label{eq:cap-truck}\\
q(s)\, y_{d,s} &\le Q_d, & d\in D,\ s\in S^d \label{eq:cap-drone}\\
q(s)\, z_{r,s} &\le Q_r, & r\in R,\ s\in S^r \label{eq:cap-robot}
\end{align}
{\footnotesize\noindent\emph{Code:} \texttt{truck\_capacity\_\{v\}},
\texttt{drone\_capacity\_\{d\}\_\{sid\}}, \texttt{robot\_capacity\_\{r\}\_\{sid\}}.\par}

\subsection{Endurance (hard)}
\label{sec:endurance}
\begin{align}
\delta(s)\, y_{d,s} &\le E_d, & d\in D,\ s\in S^d \label{eq:end-drone}\\
\delta(s)\, z_{r,s} &\le E_r, & r\in R,\ s\in S^r \label{eq:end-robot}
\end{align}
{\footnotesize\noindent\emph{Code:} \texttt{drone\_endurance\_\{d\}\_\{sid\}},
\texttt{robot\_endurance\_\{r\}\_\{sid\}}. Candidate generation
(Section~\ref{sec:sorties}) already excludes any sortie with $\delta(s) > E_d$ (or
$E_r$); these constraints restate the limit explicitly in the MILP.\par}

\subsection{Drone sortie--truck assignment and synchronization}
\label{sec:drone-sync}
For every $d\in D$, $s\in S^d$:
\begin{align}
\sum_{v\in T} h^L_{v,d,s} &= y_{d,s} \label{eq:drone-launch-assign}\\
\sum_{v\in T} h^R_{v,d,s} &= y_{d,s} \label{eq:drone-recover-assign}
\end{align}
and for every $v\in T$:
\begin{align}
h^L_{v,d,s} &\le
\begin{cases}
\eta_v, & \iota(s) = 0\\
\mathrm{vis}_{v,\iota(s)}, & \iota(s)\in C\\
0, & \text{otherwise}
\end{cases} \label{eq:drone-launch-sync}\\
h^R_{v,d,s} &\le
\begin{cases}
\eta_v, & \kappa(s) = n+1\\
\mathrm{vis}_{v,\kappa(s)}, & \kappa(s)\in C\\
0, & \text{otherwise}
\end{cases} \label{eq:drone-recover-sync}
\end{align}
{\footnotesize\noindent\emph{Code:} \texttt{assign\_selected\_drone\_sortie\_launch/recovery\_truck\_\{d\}\_\{sid\}}, \texttt{sync\_drone\_launch/recover\_*}. A selected sortie
must have exactly one launch truck and one recovery truck, and that truck must
physically be at the launch/recovery node (depot use $\eta_v$; customer nodes use
$\mathrm{vis}_{v,\cdot}$). The launch and recovery truck need not be the same
truck. The ``otherwise'' branch is a defensive bound that is never active for the
node sets produced by the data (launch/recovery nodes are always $0$, $n+1$, or in
$C$).\par}

\subsection{Physical-platform consistency}
\begin{align}
\sum_{s\in S^d} y_{d,s} &\le 1, & d\in D \label{eq:one-sortie-drone}\\
\sum_{s\in S^r} z_{r,s} &\le 1, & r\in R \label{eq:one-sortie-robot}
\end{align}
{\footnotesize\noindent\emph{Code:} \texttt{one\_sortie\_per\_drone\_\{d\}},
\texttt{one\_sortie\_per\_robot\_\{r\}}. Each physical drone/robot flies
at most one sortie, so no inter-sortie sequencing is needed.\par}

\subsection{Truck--platform pairing limits}
For every $v\in T$:
\begin{align}
\sum_{d\in D} \xi_{v,d} &\le \mu^{trk}_d\,\eta_v \label{eq:drone-truck-cap}\\
\sum_{s\in S^d}\big(h^L_{v,d,s} + h^R_{v,d,s}\big) &\le 2\,\xi_{v,d}, & d\in D \label{eq:drone-truck-upper}\\
\xi_{v,d} &\le \sum_{s\in S^d}\big(h^L_{v,d,s} + h^R_{v,d,s}\big), & d\in D \label{eq:drone-truck-lower}\\[4pt]
\sum_{r\in R} \zeta_{v,r} &\le \mu^{trk}_r\,\eta_v \label{eq:robot-truck-cap}\\
\sum_{s\in S^r}\big(g^L_{v,r,s} + g^R_{v,r,s}\big) &\le 2\,\zeta_{v,r}, & r\in R \label{eq:robot-truck-upper}\\
\zeta_{v,r} &\le \sum_{s\in S^r}\big(g^L_{v,r,s} + g^R_{v,r,s}\big), & r\in R \label{eq:robot-truck-lower}
\end{align}
{\footnotesize\noindent\emph{Code:} \texttt{max\_distinct\_drones\_assigned\_to\_truck\_\{v\}},
\texttt{drone\_truck\_pair\_upper/lower\_\{v\}\_\{d\}}, and the robot
analogues. $\xi_{v,d}=1$ iff drone $d$ launches and/or is recovered by
truck $v$ at least once (at most one launch and one recovery total, by
\eqref{eq:one-sortie-drone}, hence the coefficient $2$).\par}

\subsection{Launch/recovery node capacity limits}
For every $v\in T$ and every node $w\in V$ that is a launch node of some
$s\in S^d$:
\begin{align}
\sum_{d\in D}\sum_{\substack{s\in S^d\\ \iota(s)=w}} h^L_{v,d,s} \le
\begin{cases}
\mu^{dep}_d\,\eta_v, & w = 0\\
\mu^{trk}_d\,\mathrm{vis}_{v,w}, & w\ne 0
\end{cases} \label{eq:drone-launch-limit}
\end{align}
and for every recovery node $w\in V$ that is a recovery node of some $s\in S^d$:
\begin{align}
\sum_{d\in D}\sum_{\substack{s\in S^d\\ \kappa(s)=w}} h^R_{v,d,s} \le
\begin{cases}
\mu^{trk}_d\,\eta_v, & w = n+1\\
\mu^{trk}_d\,\mathrm{vis}_{v,w}, & w\in C
\end{cases} \label{eq:drone-recovery-limit}
\end{align}
\noindent with the symmetric pair for robots using $g^L, g^R, \mu^{dep}_r,
\mu^{trk}_r$.

{\footnotesize\noindent\emph{Code:}
\begin{itemize}[topsep=2pt,itemsep=1pt]
\item \texttt{depot\_drone\_launch\_limit\_\{v\}} / \texttt{node\_drone\_launch\_limit\_\{v\}\_\{launch\}};
\item \texttt{depot\_drone\_recovery\_limit\_\{v\}} / \texttt{node\_drone\_recovery\_limit\_\{v\}\_\{recover\}};
\item robot analogues.
\end{itemize}\par}

\subsection{Symmetry breaking}
\begin{align}
\sum_{s\in S^d} y_{d+1,s} &\le \sum_{s\in S^d} y_{d,s}, & d\in\{1,\dots,|D|-1\} \label{eq:sym-drone}\\
\sum_{s\in S^r} z_{r+1,s} &\le \sum_{s\in S^r} z_{r,s}, & r\in\{1,\dots,|R|-1\} \label{eq:sym-robot}
\end{align}
{\footnotesize\noindent\emph{Code:} \texttt{symmetry\_drone\_\{d\_prev\}\_\{d\_next\}},
\texttt{symmetry\_robot\_\{r\_prev\}\_\{r\_next\}}. Platforms of the same
type are identical and each flies at most one sortie, so platform $d+1$ may be used
only if platform $d$ is already used.\par}

\subsection{Drone and robot time windows}
\label{sec:drone-robot-tw}
\begin{align}
a^d_0 &= 0, \qquad a^r_0 = 0 \label{eq:start-time-platforms}\\
a^d_i &\ge e_i\, \sigma^d_i, & i\in C \label{eq:open-drone}\\
a^r_i &\ge e_i\, \sigma^r_i, & i\in C \label{eq:open-robot}\\
\ell^d_i &\ge a^d_i - \bar{e}_i, & i\in C \label{eq:late-drone-def}\\
\ell^r_i &\ge a^r_i - \bar{e}_i, & i\in C \label{eq:late-robot-def}\\
a^d_i &\le \bar{e}_i + \ell^d_i + M\,(1-\sigma^d_i), & i\in C \label{eq:tw-drone}\\
a^r_i &\le \bar{e}_i + \ell^r_i + M\,(1-\sigma^r_i), & i\in C \label{eq:tw-robot}
\end{align}
{\footnotesize\noindent\emph{Code:} \texttt{start\_time\_d}, \texttt{start\_time\_r},
\texttt{open\_drone\_\{i\}}, \texttt{open\_robot\_\{i\}}, \texttt{late\_drone\_def\_\{i\}},
\texttt{late\_robot\_def\_\{i\}}, \texttt{time\_window\_drone\_\{i\}},
\texttt{time\_window\_robot\_\{i\}}.\par}

\paragraph{Reporting note.} For a customer $i$ with $\sigma^d_i=0$,
\eqref{eq:open-drone} reduces to $a^d_i\ge 0$ and, since $\lambda_W$ drives
$\ell^d_i\to 0$, \eqref{eq:tw-drone} reduces to $a^d_i\le \bar{e}_i$. The variable
$a^d_i$ then has no effect on the objective, so the solver may set it to either
bound of $[0,\bar{e}_i]$ with no physical meaning. The Arrival Times sheet flags such
rows as \texttt{not\_served\_value\_arbitrary} (same for $a^r_i$).

\subsection{Ordered sortie timing}
For $d\in D$, $s\in S^d$ with sequence $\sigma(s)=(\sigma_1,\dots,\sigma_m)$,
launch node $\iota(s)$, recovery node $\kappa(s)$, and for every $v\in T$:
\begin{align}
a^d_{\sigma_1} &\ge a^t_{v,\iota(s)} + \mathrm{dt}_{\iota(s),\sigma_1} - M\,(1-h^L_{v,d,s}) \label{eq:sortie-drone-launch}\\
a^t_{v,\kappa(s)} &\ge a^d_{\sigma_m} + \theta_{\sigma_m} + \mathrm{dt}_{\sigma_m,\kappa(s)} - M\,(1-h^R_{v,d,s}) \label{eq:sortie-drone-recover}
\end{align}
and for every consecutive pair $(\sigma_l,\sigma_{l+1})$ in the sequence:
\begin{align}
a^d_{\sigma_{l+1}} &\ge a^d_{\sigma_l} + \theta_{\sigma_l} + \mathrm{dt}_{\sigma_l,\sigma_{l+1}} - M\,(1-y_{d,s}) \label{eq:sortie-drone-seq}
\end{align}
The robot analogues replace $a^d, h^L, h^R, y, \mathrm{dt}$ with $a^r, g^L, g^R, z,
\mathrm{rt}$:
\begin{align}
a^r_{\sigma_1} &\ge a^t_{v,\iota(s)} + \mathrm{rt}_{\iota(s),\sigma_1} - M\,(1-g^L_{v,r,s}) \label{eq:sortie-robot-launch}\\
a^t_{v,\kappa(s)} &\ge a^r_{\sigma_m} + \theta_{\sigma_m} + \mathrm{rt}_{\sigma_m,\kappa(s)} - M\,(1-g^R_{v,r,s}) \label{eq:sortie-robot-recover}\\
a^r_{\sigma_{l+1}} &\ge a^r_{\sigma_l} + \theta_{\sigma_l} + \mathrm{rt}_{\sigma_l,\sigma_{l+1}} - M\,(1-z_{r,s}) \label{eq:sortie-robot-seq}
\end{align}
{\footnotesize\noindent\emph{Code:} \texttt{sortie\_time\_drone\_launch/recover/seq\_*};
\texttt{sortie\_time\_robot\_launch/recover/seq\_*}. If a sortie $s$ is
selected ($y_{d,s}=1$ or $z_{r,s}=1$), the platform's timeline runs from the launch
truck's arrival at $\iota(s)$, through the ordered sequence with service times,
to no later than the recovery truck's arrival at $\kappa(s)$.\par}

\subsection{Truck routing completeness}
\label{sec:routing}
For every $v\in T$:
\begin{align}
\sum_{j:(0,j)\in A} x_{v,0,j} &= \eta_v \label{eq:depot-depart}\\
\sum_{i:(i,n+1)\in A} x_{v,i,n+1} &= \eta_v \label{eq:depot-return}
\end{align}
For every $v\in T$ and $\mathrm{node}\in C$ (flow conservation):
\begin{align}
\sum_{i:(i,\mathrm{node})\in A} x_{v,i,\mathrm{node}} = \sum_{j:(\mathrm{node},j)\in A} x_{v,\mathrm{node},j} \label{eq:flow}
\end{align}
For every $v\in T,\ i\in C$:
\begin{align}
\mathrm{vis}_{v,i} &= \sum_{h:(h,i)\in A} x_{v,h,i} \label{eq:visit-def}
\end{align}
For every $i \in C$:
\begin{align}
\sigma^t_i &= \sum_{v\in T} \mathrm{vis}_{v,i} \label{eq:truck-service-restated}
\end{align}
For every $v\in T$:
\begin{align}
\sum_{i\in C} \mathrm{vis}_{v,i} &\le n\,\eta_v \label{eq:used-upper}\\
\eta_v &\le \sum_{i\in C} \mathrm{vis}_{v,i} \label{eq:used-lower}\\
x_{v,i,j} &\le \eta_v, & (i,j)\in A \label{eq:arc-activation}
\end{align}
{\footnotesize\noindent\emph{Code:}
\begin{itemize}[topsep=2pt,itemsep=1pt]
\item \texttt{op\_depot\_departure\_\{v\}} / \texttt{op\_depot\_return\_\{v\}};
\item \texttt{op\_flow\_\{v\}\_\{node\}};
\item \texttt{op\_visit\_inbound\_\{v\}\_\{i\}};
\item \texttt{op\_truck\_service\_visit\_\{i\}};
\item \texttt{op\_used\_visit\_activation\_\{v\}} / \texttt{op\_used\_requires\_service\_\{v\}} / \texttt{op\_arc\_activation\_\{v\}\_\{i\}\_\{j\}}.
\end{itemize}
These constraints are not written explicitly in the original PDF, which
defines $x_{v,i,j}$ but not full route-continuity; they are the standard
constraints needed for $x$ to form valid physical truck routes.\par}

\subsection{Truck timing and time windows}
\begin{align}
a^t_{v,0} &= 0, & v\in T \label{eq:truck-start-time}\\
a^t_{v,j} &\ge a^t_{v,i} + \theta_i\,[i\in C] + \mathrm{tt}_{ij} - M\,(1-x_{v,i,j}), & v\in T,\ (i,j)\in A \label{eq:truck-time-propagation}\\
a^t_{v,i} &\ge e_i - M\,(1-\mathrm{vis}_{v,i}), & v\in T,\ i\in C \label{eq:truck-open}\\
\ell^t_i &\ge a^t_{v,i} - \bar{e}_i - M\,(1-\mathrm{vis}_{v,i}), & v\in T,\ i\in C \label{eq:truck-late}
\end{align}
{\footnotesize\noindent\emph{Code:} \texttt{op\_start\_time\_truck\_\{v\}},
\texttt{op\_time\_propagation\_\{v\}\_\{i\}\_\{j\}},
\texttt{op\_open\_truck\_\{v\}\_\{i\}} / \texttt{op\_late\_truck\_\{v\}\_\{i\}}. Here $[i\in C]$ is $1$ if $i\in C$ (so $\theta_i$ is added) and $0$ if
$i\in P$ (depots have no service time). Every used truck departs the start depot at
clock time $0$.\par}

\subsection{Route-duration constraints}
For every $v\in T$:
\begin{align}
\rho^t_v &\ge a^t_{v,n+1} - T_{\max} - M\,(1-\eta_v) \label{eq:route-late-def}\\
a^t_{v,n+1} &\le T_{\max} + \rho^t_v + M\,(1-\eta_v) \label{eq:route-late-bound}\\
\rho^t_v &\le M\,\eta_v \label{eq:route-late-activation}\\
\rho^t_v &\le \rho^* + M\,(1-\eta_v) \label{eq:route-late-global-bound}
\end{align}
{\footnotesize\noindent\emph{Code:} \texttt{op\_route\_duration\_late\_truck\_\{v\}},
\texttt{op\_route\_duration\_truck\_\{v\}},
\texttt{op\_route\_duration\_late\_truck\_activation\_\{v\}},
\texttt{op\_route\_duration\_late\_truck\_global\_bound\_\{v\}}.
Because every truck departs at time $0$ (\eqref{eq:truck-start-time}), $a^t_{v,n+1}$
\emph{is} that truck's route duration. $\rho^*$ is forced, by
\eqref{eq:route-late-activation}--\eqref{eq:route-late-global-bound} together with
the objective minimizing $\lambda_T\rho^*$, to equal $\max_{v\in T,\,\eta_v=1}
\rho^t_v$, i.e. the worst per-truck route-duration excess over $T_{\max}$.\par}

\paragraph{Relation to drone/robot routes.} The original PDF writes the
route-duration condition for $m\in\{t,d,r\}$ (trucks, drones, robots
separately). This implementation enforces it only through the truck end-depot
arrival $a^t_{v,n+1}$. A drone/robot finishing a sortie late forces its recovery
truck to wait (via \eqref{eq:sortie-drone-recover}/\eqref{eq:sortie-robot-recover}),
which propagates into that truck's $a^t_{v,n+1}$ and hence into $\rho^t_v$ and
$\rho^*$. So drone/robot lateness is captured indirectly through the truck term,
not as a separate $a^d_{n+1}, a^r_{n+1}$ penalty. See Section~\ref{sec:deviations}.

\subsection{MTZ subtour elimination}
For every $v\in T$:
\begin{align}
u_{v,i} &\le n\,\mathrm{vis}_{v,i}, & i\in C \label{eq:mtz-upper}\\
u_{v,i} &\ge \mathrm{vis}_{v,i}, & i\in C \label{eq:mtz-lower}\\
u_{v,i} - u_{v,j} + n\,x_{v,i,j} &\le n-1, & i\ne j\in C,\ (i,j)\in A \label{eq:mtz}
\end{align}
{\footnotesize\noindent\emph{Code:} \texttt{op\_mtz\_visit\_upper/lower\_\{v\}\_\{i\}},
\texttt{op\_mtz\_\{v\}\_\{i\}\_\{j\}}. Standard Miller--Tucker--Zemlin
constraints; not in the PDF, needed to forbid customer-only subtours disconnected
from the depot.\par}

\section{Lossless Fleet Trimming}
\label{sec:trimming}

Before $D$ and $R$ are fixed, the raw fleet sizes from \texttt{parameters.json}
(\texttt{NUM\_DRONES}, \texttt{NUM\_ROBOTS}) are trimmed to
\begin{align}
|D| &= \max\big(1,\ \min(\mu^{trk}_d\cdot|T|,\ n)\big) \label{eq:trim-drone}\\
|R| &= \max\big(1,\ \min(\mu^{trk}_r\cdot|T|,\ n)\big) \label{eq:trim-robot}
\end{align}
\noindent\emph{Code:} the fleet-trimming block. This is lossless: \eqref{eq:drone-truck-cap}
already caps usable drones at $\mu^{trk}_d\cdot|T|$ per type, and
\eqref{eq:one-sortie-drone} plus \eqref{eq:service} cap usable drones at $n$ (one
sortie serves at least one customer, each customer served once), so any platform
beyond $|D|$ (or $|R|$) is provably unused in every feasible solution.

\section{Deviations from the Original PDF Formulation}
\label{sec:deviations}

\begin{itemize}
\item \textbf{Endurance and capacity are hard constraints.} The PDF's Eq.~12
penalizes endurance/capacity violations with $\lambda^d_E, \lambda^r_E, \lambda_Q$
to allow \emph{heuristic exploration of infeasible solutions}. This MILP instead
enforces \eqref{eq:cap-truck}--\eqref{eq:end-robot} as hard constraints (and
additionally at sortie-generation time), so $\lambda^d_E, \lambda^r_E, \lambda_Q$
are unused.
\item \textbf{Route-duration penalty is truck-indexed only}, as discussed after
\eqref{eq:route-late-global-bound}, rather than summed separately over
$m\in\{t,d,r\}$ as in the PDF's Eq.~8/Eq.~12. Drone/robot lateness still affects
$\rho^*$ indirectly through synchronization.
\item \textbf{Big-M, waiting penalty, and raw cost fields}
(\texttt{big\_M}, \texttt{waiting\_penalty\_weight}, \texttt{truck\_variable\_cost},
\texttt{drone\_departure\_cost}) from the raw TDRP-TW parameters are superseded:
$M$ is computed internally (Section~\ref{sec:parameters}, parameter table), platform waiting is
reported via the Sortie Timing Audit rather than penalized, and travel costs use
$C_{veh}/C_{drone}/C_{rob}$ with the fixed-cost terms \eqref{eq:fixcost-truck}--\eqref{eq:fixcost-robot}.
\item \textbf{Operational routing constraints}
(\eqref{eq:depot-depart}--\eqref{eq:mtz}) are added because the PDF defines
$x_{v,i,j}$ but does not write full route-continuity constraints; without them
$x$ does not necessarily form a connected truck route from depot to depot.
\item \textbf{Makespan} is computed after solving from $a^t_{v,n+1}$ and platform
finish times for reporting only. It is not a decision variable and is not part of
$Z$, matching the PDF's objective structure $P_{\text{PDF}} = Z + \text{penalty terms}$
(Eq.~12, where $P_{\text{PDF}}$ denotes the PDF's total-cost objective, not the depot
set $P$ of Section~\ref{sec:sets}), which
also does not include makespan in the objective.
\end{itemize}

A full parameter-by-parameter audit is in \texttt{parameter\_usage\_audit.csv} and
the README's ``Parameter Usage Audit'' section.

\end{document}
